Ce travail pratique s'inscrit dans le domaine de l'optique géométrique, 
et plus particulièrement dans l'étude du comportement de la lumière à travers un prisme.
L'objectif du T.P. est double. 
D'abord, il s'agit de mettre en pratique le maniement d'un goniomètre. 
Enfin, il s'agit également de mesurer expérimentalement les caractéristiques du prisme, 
notamment son minimum de déviation et son indice de réfraction.
La problématique posée est la suivante : 
comment les mesures de l'angle utile et de l'angle de minimum de déviation permettent-elles de déterminer avec précision l'indice de réfraction du verre constituant le prisme ?